\documentclass[paper=a4, fontsize=11pt]{scrartcl}

\usepackage[T1]{fontenc}
\usepackage{fourier}
\usepackage[english]{babel}
\usepackage{amsmath,amsfonts,amsthm}

\usepackage{sectsty}
\allsectionsfont{\centering \normalfont\scshape}

\usepackage{fancyhdr}
\pagestyle{fancyplain}
\fancyhead{}
\fancyfoot[L]{}
\fancyfoot[C]{}
\fancyfoot[R]{\thepage}
\renewcommand{\headrulewidth}{0pt}
\renewcommand{\footrulewidth}{0pt}
\setlength{\headheight}{13.6pt}

\setlength\parindent{0pt}

% Title helpers
\newcommand{\horrule}[1]{\rule{\linewidth}{#1}}

\title{
\normalfont \normalsize
\textsc{Artificial Intelligence - 4,906,1.00 - Sheet 3 Spring 2026} \\[10pt]
\horrule{0.5pt} \\[0.4cm]
\huge Solutions \\[-2pt]
\horrule{2pt} \\[0.5cm]
}

\author{
\textbf{Group 12}\\
Valeria Pointet\\
Keanu Belo da Silva
}

\date{\normalsize\today}

\begin{document}
\maketitle

% -------------------- 1 --------------------
\section{Measuring Performance}

\subsection*{Question 1.1}
Precision measures how many predicted positives are actually correct.

\[
\text{Precision} = \frac{TP}{TP + FP}
\]

Recall measures how many actual positives were found.

\[
\text{Recall} = \frac{TP}{TP + FN}
\]

Higher precision is preferable when false positives are costly.  
Higher recall is preferable when missing positives is costly.

A metric that combines both is the F1-score.

\[
F1 = \frac{2 \cdot \text{Precision} \cdot \text{Recall}}{\text{Precision} + \text{Recall}}
\]

\subsection*{Question 1.2}
Top-k accuracy means that a prediction is counted as correct if the true label is among the top k predicted labels.

Top-5 accuracy is often preferred when:
\begin{itemize}
    \item there are many classes
    \item several classes are similar
    \item one exact prediction is too strict
\end{itemize}

\subsection*{Question 1.3}
The system retrieves 6 relevant documents out of 10 retrieved documents.

\[
\text{Precision} = \frac{6}{10} = 0.6
\]

There are 15 relevant documents in total.

\[
\text{Recall} = \frac{6}{15} = 0.4
\]

So the F1-score is:

\[
F1 = \frac{2 \cdot 0.6 \cdot 0.4}{0.6 + 0.4} = 0.48
\]

So the F1-score is:

\[
0.48
\]

\subsection*{Question 1.4}
The total number of samples is:

\[
40+5+5+6+30+4+3+7+25 = 125
\]

The correctly classified samples are the diagonal entries:

\[
40+30+25 = 95
\]

So the overall accuracy is:

\[
\frac{95}{125} = 0.76
\]

Now we compute recall for each class.

For class A:

\[
\frac{40}{40+5+5} = \frac{40}{50} = 0.8
\]

For class B:

\[
\frac{30}{6+30+4} = \frac{30}{40} = 0.75
\]

For class C:

\[
\frac{25}{3+7+25} = \frac{25}{35} \approx 0.71
\]

So class A has the highest recall.

% -------------------- 2 --------------------
\section{Confidence-based Benchmarking}

\subsection*{Question 2.1}
A sample is classified as positive if its score is greater than or equal to the threshold.

There are 3 positive samples and 3 negative samples.

\textbf{Threshold $t=1.0$}

No sample is classified as positive.

\[
TP=0,\quad FP=0,\quad TN=3,\quad FN=3
\]

\[
TPR=\frac{0}{3}=0,\quad FPR=\frac{0}{3}=0
\]

\textbf{Threshold $t=0.8$}

Positive samples: 1, 2

\[
TP=1,\quad FP=1,\quad TN=2,\quad FN=2
\]

\[
TPR=\frac{1}{3},\quad FPR=\frac{1}{3}
\]

\textbf{Threshold $t=0.6$}

Positive samples: 1, 2, 3, 4

\[
TP=2,\quad FP=2,\quad TN=1,\quad FN=1
\]

\[
TPR=\frac{2}{3},\quad FPR=\frac{2}{3}
\]

\textbf{Threshold $t=0.4$}

Positive samples: 1, 2, 3, 4, 5

\[
TP=3,\quad FP=2,\quad TN=1,\quad FN=0
\]

\[
TPR=1,\quad FPR=\frac{2}{3}
\]

\textbf{Threshold $t=0.3$}

All samples are classified as positive.

\[
TP=3,\quad FP=3,\quad TN=0,\quad FN=0
\]

\[
TPR=1,\quad FPR=1
\]

So the table is:

\[
\begin{array}{c|c|c}
\text{Threshold} & \text{TPR} & \text{FPR} \\
\hline
1.0 & 0 & 0 \\
0.8 & \frac{1}{3} & \frac{1}{3} \\
0.6 & \frac{2}{3} & \frac{2}{3} \\
0.4 & 1 & \frac{2}{3} \\
0.3 & 1 & 1
\end{array}
\]

The ROC points are:

\[
(0,0),\left(\frac{1}{3},\frac{1}{3}\right),\left(\frac{2}{3},\frac{2}{3}\right),\left(\frac{2}{3},1\right),(1,1)
\]

The point $(0,1)$ is the ideal point on the ROC curve.  
It means no false positives and all positives found.

\subsection*{Question 2.2}
PR curves are more adequate for imbalanced classification problems.

This is because ROC curves can still look good when there are many true negatives.

PR curves focus on precision and recall.  
So they show better how well the model finds the positive class.

A typical example is fraud detection.

\end{document}